This survey targets a single phenomenon: large language models that produce fluent, confident, and yet incorrect text. We unify the literature around six anchors. The first is scope: text, vision-language, audio, and embodied models, with primary focus on text-only LLMs from GPT-3 (2020) through DeepSeek-R1 (January 2025), GPT-4o (May 2024), Claude 3.5 Sonnet (June 2024), Gemini 1.5 Pro (February 2024), Gemini 2.5 (2025), and LLaMA-3-405B (April 2024). The second is taxonomy: a four-axis decomposition by cause (data, model, training, inference), surface form (entity, relation, numerical, citation, code), modality (text, vision-language, audio, embodied), and mitigation family (retrieval, decoding, refinement, alignment, editing). The third is the historical arc, which runs from neural-machine-translation mistranslations in 2018, through the intrinsic/extrinsic split of Maynez et al.~(ACL 2020), through TruthfulQA in 2021 by Lin, Hilton, and Evans, through the ChatGPT and GPT-4 inflection of late 2022 and early 2023, to the consolidation phase of 2024--2026 marked by the \emph{Nature} paper of Farquhar et al.~on semantic entropy (vol.~630, June 2024). The fourth anchor is the set of key methods. SelfCheckGPT (Manakul et al.~2023) reaches AUROC 0.92 on WikiBio-GPT3. DoLa (Chuang et al.~2023) lifts TruthfulQA-MC1 by 6.8 points. FActScore (Min et al.~2023) introduces atomic-precision factuality. Chain-of-Verification (Dhuliawala et al.~2024) gains 29 points on Wikidata-Listed. Infer\ldots{}
